% =====================================================================
%  Mise en forme imposee par « Rapport Stage Pedagogique.pdf »
% =====================================================================
%  Police Times 12 pt, marges 2,5 cm, interligne 1,15, justification,
%  alinea de premiere ligne 0,50 cm, pagination en bas a droite sous la
%  forme n/N. Chaque valeur ci-dessous est reprise du guide.
% =====================================================================

\setstretch{1.15}
\setlength{\parindent}{0.5cm}
\setlength{\parskip}{0.15\baselineskip}

% ---------------------------------------------------------------- Titres
\definecolor{bleuTitre}{HTML}{1F3864}
\definecolor{grisFilet}{HTML}{9AA5B1}
\definecolor{grisFond}{HTML}{F2F4F7}

\titleformat{\chapter}[display]
  {\normalfont\sffamily\bfseries\color{bleuTitre}}
  {\filright\large\MakeUppercase{\chaptertitlename}\ \Huge\thechapter}
  {8pt}{\huge\filright}
\titlespacing*{\chapter}{0pt}{-20pt}{28pt}

\titleformat{\section}{\normalfont\sffamily\large\bfseries\color{bleuTitre}}
  {\thesection}{0.6em}{}
\titleformat{\subsection}{\normalfont\sffamily\normalsize\bfseries\color{bleuTitre}}
  {\thesubsection}{0.6em}{}
\titleformat{\subsubsection}{\normalfont\sffamily\small\bfseries\color{bleuTitre}}
  {\thesubsubsection}{0.6em}{}

\titlespacing*{\section}      {0pt}{14pt}{6pt}
\titlespacing*{\subsection}   {0pt}{11pt}{4pt}
\titlespacing*{\subsubsection}{0pt}{9pt}{3pt}

\setcounter{secnumdepth}{3}
\setcounter{tocdepth}{2}

% ---------------------------------------------------------------- En-tetes
\pagestyle{fancy}
\fancyhf{}
\renewcommand{\headrulewidth}{0.4pt}
\renewcommand{\footrulewidth}{0pt}
\fancyhead[L]{\small\sffamily\nouppercase{\leftmark}}
\fancyfoot[R]{\small\sffamily\thepage/\pageref{LastPage}}

\fancypagestyle{plain}{%
  \fancyhf{}%
  \renewcommand{\headrulewidth}{0pt}%
  \fancyfoot[R]{\small\sffamily\thepage/\pageref{LastPage}}%
}

\renewcommand{\chaptermark}[1]{\markboth{\thechapter.\ #1}{}}

% ---------------------------------------------------------------- Legendes
\captionsetup[figure]{name=Figure}
\captionsetup[table]{name=Tableau}
\renewcommand{\cftfigpresnum}{Figure~}
\renewcommand{\cfttabpresnum}{Tableau~}
\setlength{\cftfignumwidth}{4.2em}
\setlength{\cfttabnumwidth}{4.6em}

% Les tableaux portent leur legende au-dessus, les figures en dessous.
\captionsetup[table]{position=above,skip=4pt}
\captionsetup[figure]{position=below,skip=6pt}

% ---------------------------------------------------- Pages liminaires
%  POURQUOI DES PAGES BLANCHES APPARAISSENT
%  La classe est chargee en « twoside, openright » : chaque nouvelle
%  piece commence sur une page de DROITE, comme il est d'usage pour un
%  document imprime en recto-verso. Quand la piece precedente se termine
%  elle-meme sur une page de droite, LaTeX intercale une page blanche.
%  C'est voulu, et c'est ce que demande le guide, qui recommande
%  l'impression recto-verso.
%
%  Les CHAPITRES conservent ce comportement : commencer un chapitre au
%  dos d'un autre serait fautif.
%
%  Les pages liminaires, elles, n'ont pas cette exigence. Le reglage
%  ci-dessous permet de choisir :
%      \clearpage        -> pas de page blanche entre elles  (defaut)
%      \cleardoublepage  -> chacune sur une page de droite
\newcommand{\sautLiminaire}{\clearpage}

% ---------------------------------------------------------------- Divers
\renewcommand{\arraystretch}{1.25}
\setlist{itemsep=2pt,parsep=0pt,topsep=4pt}
\clubpenalty=10000      % pas de ligne isolee en bas de page
\widowpenalty=10000     % pas de ligne isolee en haut de page
\raggedbottom
